\subsection{One-use privacy in the shared byte random oracle}
\label{sec:native-random-oracle}

The public oracle accepts arbitrary byte strings and returns independent
uniform 32-byte digests on previously unseen inputs, caching each answer.
The experiment starts with an empty cache. Before releasing any view, the
constructor samples independent uniform Boolean vectors
\(\mu,k\in\{0,1\}^n\) and an independent uniform 32-byte nonce \(r\).
Here \(\varepsilon(0)=-1\), \(\varepsilon(1)=1\) coordinatewise, and
\[
 B_k=A\operatorname{diag}(\varepsilon(k)),\qquad
 z=d\varepsilon(\mu),\qquad u=B_k\varepsilon(\mu).
\]
The source opening \((\rho,S)\), exact opening \(a\), and public matrix \(A\)
are fixed independently of these samples. Choose a common descriptor with
\(\ell=|\rho|\) and a coordinate bound \(K\) satisfying
\[
\begin{gathered}
 |S_{ij}|\leq K,\quad |A_{ij}|\leq K,\quad |d|\leq K,\quad |a|\leq K,\\
 \sum_j |A_{ij}|\leq K\qquad\text{for every row }i.
\end{gathered}
\]
Every generated private carrier then satisfies the same bounded relation:
the target entries have magnitude \(|A_{ij}|\), the message entries have
magnitude \(|d|\), and the coexact entries are bounded by the row sums.

The constructor first obtains the operator root by querying this same byte
oracle at the complete canonical preimage for
\(\operatorname{Com}(B_k;r)\). It then queries the source-credential bytes
and full ciphertext bytes, and releases
\(\operatorname{encodeInstance}(x_w)\) together with that ciphertext.
The decoder receives only this view, uniform random sampling and the public
byte oracle. It has no private-derivation interface or earlier oracle
transcript. Its query budget \(Q\) counts all of its postrelease byte-oracle
queries, on every input view and every execution path.

The proof uses the hidden-seed simulation theorem of the pinned VCVio
library \cite{vcvio2026library}. The real handler shares one cache between
public queries and the constructor's private operator derivation; the ideal
handler uses an independent derivation table. Injectivity of the canonical
packet/nonce encoding implies that any fixed byte address is an operator
opening for at most one nonce seed. A uniform 32-byte seed therefore hits
that address with probability at most \(256^{-32}\). The adaptive event
comparison charges this term once per public byte-oracle query.

\begin{theorem}[Native recovery in the random-oracle model]
\label{thm:coupled-native-rom}
Fix \(A\in\mathbb Z^{m\times n}\) with \(A^{\mathsf T}A=dI\), a fixed
source credential opening and exact opening \(a\). In the fresh-cache
experiment above, every decoder making at most \(Q\) public byte-oracle
queries satisfies
\[
 \Pr_{\mathrm{RO}}[\text{recover the complete Boolean message}]
 \leq 2^{-n}+\frac{Q+2}{256^{32}}.
\]
If \(d>0\), the same upper bound holds for a partial extractor of the
original integer harmonic coefficients \(d\varepsilon(\mu)\). The two additional
queries are the constructor's source-credential and full-ciphertext digest
calls.
\end{theorem}

\begin{proof}
The constructor makes exactly one private derivation, before any public
calls, and then two public digest calls. Uniform draws are uncharged. The
lifted decoder has no private-derivation access, so the complete program
has public-query budget \(Q+2\). Hidden-seed event comparison bounds real
recovery by ideal recovery plus \((Q+2)256^{-32}\).

In the ideal game, the first derivation returns an independent uniform
operator root. All later calls use only the public byte oracle. Projecting
away the separate derivation table leaves that same public-oracle execution.
Reindex the fresh key by the fixed-message mask
\(\varepsilon(k)\mathbin\odot\varepsilon(\mu)\). This is a permutation of the uniform
key space, and the native normal form makes the ciphertext depend only on
this mask. The ideal operator root, ciphertext, two public digests and
entire subsequent adaptive decoder execution are consequently independent
of the uniform message. Complete-message guessing has mass exactly
\(2^{-n}\). For a partial coefficient extractor, the same factorization and
positive scaling give success probability at most \(2^{-n}\); failure and
invalid integer outputs cannot increase that event.
\end{proof}

\begin{theorem}[One-use chosen-message indistinguishability]
\label{thm:coupled-native-rom-challenge}
Fix any pair of Boolean messages before the experiment. Replace the uniform
message draw above by a fair challenge bit \(b\) and encrypt the \(b\)-th
message. The key, nonce, initial cache and released view are otherwise as
above. A public adversary making at most \(Q\) adaptive byte-oracle queries
after release has
\[
 \left|\Pr_{\mathrm{RO}}[\widehat b=b]-\tfrac12\right|
 \leq \frac{Q+2}{256^{32}}.
\]
\end{theorem}

\begin{proof}
In the ideal experiment the operator root is independent and the uniform
key can be reindexed by the mask for either chosen message. The complete
public execution, including the adaptive query answers, therefore has the
same law for \(b=0\) and \(b=1\). A Boolean guess is correct with probability
\(\tfrac12\). Apply the same hidden-seed event comparison to a correct
challenge-bit guess and then to its complementary event. The query budget
remains \(Q+2\).
\end{proof}

\begin{corollary}[Native recovery with deterministic key disclosure]
\label{cor:coupled-native-rom-disclosure}
In the experiment of \cref{thm:coupled-native-rom}, also give the decoder
\(L(k)\) for a fixed deterministic function \(L\) of the sampled sign key.
The disclosure has no nonce input and makes no oracle queries. Let
\(M=|\operatorname{im} L|\). A decoder making at most \(Q\) public byte-oracle
queries on every execution path has
\[
 \Pr_{\mathrm{RO}}[\text{recover the target}]
 \leq \frac{M}{2^n}+\frac{Q+2}{256^{32}}
\]
for any injective encoding of the complete message. This includes the
Boolean message and, for \(d>0\), the original integer harmonic coefficients
with a possible failure output. If \(L\) discloses \(s\) bits, the first term
is at most \(2^{s-n}\).
\end{corollary}

\begin{proof}
The query budget remains \(Q+2\), including the two honest digest calls.
The same hidden-seed comparison adds \((Q+2)256^{-32}\) to the ideal
recovery bound. In the ideal game the mask is uniform independently of the
key, since it permutes the uniform message for each key. At a fixed mask
and disclosed value, the decoder has one output distribution, including
its adaptive public queries. The compatible keys in that disclosure fiber
give distinct targets, so the probabilities assigned to those targets
sum to at most one. Each key has weight \(2^{-n}\) at this mask, hence the
fiber contributes at most \(2^{-n}\). Sum over the \(M\) disclosed values
and average over the mask.
\end{proof}

The bounds average over one shared random oracle in the one-use
instance/ciphertext experiment. The opening nonce supplies the
\(256^{32}=2^{256}\) address-guessing denominator. The payload length \(n\)
sets the uniform-message guessing baseline. Concrete hash instantiation,
prechallenge queries, repeated oracle challenges, and additional public
proof or storage transcripts require corresponding security experiments.
The key-reuse result above already applies to the commitment-bearing view
because the ciphertext bytes are included in it.
